\section{Results}
\label{sec:results}

All results are five-fold mean best test accuracy \(\pm\) fold-level standard
deviation unless stated otherwise.

\subsection{Overall Benchmark}
\label{sec:res-overall}

The main benchmark asks whether residual topology has a stable winner across
datasets and operators. It does not. Across the 24 dataset--operator
combinations (Table~\ref{tab:operator_wins}), MatrixRes is the most frequent
winner with 9 wins and the best mean rank (2.12); MatrixResGated and Plain win 5
combinations each, HorizontalRes wins 3, and VerticalRes wins 2. Matrix-style
reuse therefore has the most favorable aggregate profile, but every topology is
strongest in some setting. The per-dataset winners are mixed: the matrix family
dominates MUTAG, AIDS, and Mutagenicity, whereas PROTEINS alternates between
Plain and HorizontalRes and DD alternates among HorizontalRes, MatrixRes, and
VerticalRes. ENZYMES remains difficult and shows no stable preference.

\begin{table}[!t]
\centering
\caption{Winner Counts and Mean Ranks over the 24 Dataset--Operator
Combinations}
\label{tab:operator_wins}
\footnotesize
\begin{tabular}{@{}lccc@{}}
\toprule
Model & Wins & Mean rank & \multicolumn{1}{l}{GCNConv wins} \\
\midrule
Plain          & 5 & 3.04 & PROTEINS \\
VerticalRes    & 2 & 3.38 & --- \\
HorizontalRes  & 3 & 3.54 & DD \\
MatrixRes      & \textbf{9} & \textbf{2.12} & MUTAG, AIDS, Mutagenicity \\
MatrixResGated & 5 & 2.92 & ENZYMES \\
\bottomrule
\end{tabular}
\end{table}

ENZYMES should be read cautiously, because the absolute accuracies are low for
all models. Excluding it leaves MatrixRes with 9 wins over 20 combinations and a
mean rank of 2.05, so the aggregate preference is unchanged.

To remove variation caused by the operator, we also report the GCNConv slice
separately (Table~\ref{tab:main_benchmark}, Fig.~\ref{fig:main_benchmark}). The
same conclusion holds: MatrixRes is the strongest choice on MUTAG, AIDS, and
Mutagenicity, while simpler residual paths remain preferable on PROTEINS, DD,
and ENZYMES. Residual topology matters, but its value depends on the dataset.

\begin{table}[!t]
\centering
\caption{GCNConv Benchmark at \(B=3\) (Mean \(\pm\) Std. over Five Folds)}
\label{tab:main_benchmark}
\footnotesize
\begin{tabular}{@{}lccccc@{}}
\toprule
Dataset & Plain & Vertical & Horizontal & Matrix & MatrixGated \\
\midrule
PROTEINS     & \(\mathbf{0.7044}\) & 0.6993 & 0.6992 & 0.6973 & 0.6886 \\
DD           & 0.7053 & 0.7143 & \(\mathbf{0.7181}\) & 0.7127 & 0.7141 \\
ENZYMES      & 0.2683 & 0.2550 & 0.2633 & 0.2750 & \(\mathbf{0.2933}\) \\
MUTAG        & 0.7556 & 0.7451 & 0.7343 & \(\mathbf{0.7609}\) & 0.7556 \\
AIDS         & 0.8215 & 0.8210 & 0.8130 & \(\mathbf{0.8365}\) & 0.8150 \\
Mutagenicity & 0.7644 & 0.7667 & 0.7738 & \(\mathbf{0.7784}\) & 0.7701 \\
\bottomrule
\end{tabular}
\end{table}

\begin{figure}[!t]
\centering
\includegraphics[width=0.98\columnwidth]{figures/exp/fig_main_benchmark_gcnconv.pdf}
\caption{GCNConv benchmark at \(B=3\). Fold-level standard deviations range
from 0.005 to 0.072 across cells.}
\label{fig:main_benchmark}
\end{figure}

\subsection{Synthetic Multi-Class Scaling}
\label{sec:res-synthetic}

The synthetic benchmark asks whether residual topology behaves differently as
structural class granularity increases under controlled graph-generation rules.
Table~\ref{tab:synthetic_class_scaling} summarizes the trend averaged over the
five families and four operators. Low-class settings do not favour matrix-style
reuse: Plain is best at \(C=2\) and HorizontalRes at \(C=3\). From \(C=4\)
onward matrix-style reuse becomes more competitive, and MatrixResGated is the
strongest mean-accuracy model at \(C=4\), \(C=5\), \(C=6\), and \(C=8\). Over the
higher-class regime \(C=4,\ldots,8\), MatrixResGated reaches 0.6943 accuracy and
0.6324 normalized accuracy against 0.6827 and 0.6183 for Plain, while MatrixRes
has the best macro-F1 (0.6880). MatrixRes also degrades least from \(C=2\) to
\(C=8\) (0.3479 points, against 0.3720 for Plain), suggesting that dense local
matrix reuse is relatively stable as the label space becomes finer.

\begin{table}[!t]
\centering
\caption{Synthetic Scaling Averaged over Five Graph Families and Four
Operators}
\label{tab:synthetic_class_scaling}
\footnotesize
\begin{tabular}{@{}clcccc@{}}
\toprule
Classes & Best model & Plain & Matrix & MatrixGated & Best NAcc \\
\midrule
2 & Plain          & \(\mathbf{0.9585}\) & 0.9481 & 0.9578 & 0.9170 \\
3 & HorizontalRes  & 0.8774 & 0.8600 & 0.8772 & 0.8200 \\
4 & MatrixResGated & 0.8002 & 0.8083 & \(\mathbf{0.8111}\) & 0.7481 \\
5 & MatrixResGated & 0.7233 & 0.7280 & \(\mathbf{0.7296}\) & 0.6621 \\
6 & MatrixResGated & 0.6719 & 0.6827 & \(\mathbf{0.6878}\) & 0.6254 \\
7 & VerticalRes    & 0.6314 & 0.6429 & 0.6395 & 0.5945 \\
8 & MatrixResGated & 0.5865 & 0.6002 & \(\mathbf{0.6035}\) & 0.5469 \\
\bottomrule
\end{tabular}
\end{table}

\begin{figure}[!t]
\centering
\includegraphics[width=0.98\columnwidth]{figures/exp/fig_synthetic_multiclass_scaling.pdf}
\caption{Synthetic scaling from two to eight classes, averaged over five graph
families and four operators.}
\label{fig:synthetic_multiclass_scaling}
\end{figure}

At the family level, MatrixRes is strongest in the higher-class BA, ER, and
regular-degree settings, where labels are tied to degree growth, edge
probability, or degree-uniform structural change. These families share a
pattern: the label space is formed by related structural regimes rather than
unrelated categories, so local matrix reuse can combine depth-wise refinement
with neighboring-branch exchange. SBM is community-driven and less favourable
to matrix-style reuse, though MatrixRes remains close to the best result.

\subsection{Branch-Count Ablation and Mechanism Analysis}
\label{sec:res-branch}

Branch count is not merely a capacity parameter, because it also changes
residual-path density. Table~\ref{tab:branch_best} reports the best and worst
branch counts per topology. On PROTEINS the useful region is moderate
(HorizontalRes peaks at \(B=4\), MatrixRes at \(B=6\), MatrixResGated at
\(B=3\)); on DD it is smaller (HorizontalRes and MatrixRes peak at \(B=2\),
MatrixResGated is strongest already at \(B=1\)). Branch expansion is therefore
not monotonic and must be matched to task structure.

\begin{table}[!t]
\centering
\caption{Branch-Count Ablation: Best and Worst Settings}
\label{tab:branch_best}
\footnotesize
\begin{tabular}{@{}llcccc@{}}
\toprule
Dataset & Model & Best \(B\) & Best acc. & Worst \(B\) & Worst acc. \\
\midrule
PROTEINS & HorizontalRes  & 4 & \(\mathbf{0.7125}\) & 1 & 0.6801 \\
PROTEINS & MatrixRes      & 6 & \(\mathbf{0.7062}\) & 1 & 0.6909 \\
PROTEINS & MatrixResGated & 3 & \(\mathbf{0.7098}\) & 4 & 0.6711 \\
DD       & HorizontalRes  & 2 & \(\mathbf{0.7275}\) & 4 & 0.7071 \\
DD       & MatrixRes      & 2 & \(\mathbf{0.7206}\) & 7 & 0.7071 \\
DD       & MatrixResGated & 1 & \(\mathbf{0.7283}\) & 7 & 0.7053 \\
\bottomrule
\end{tabular}
\end{table}

\begin{figure}[!t]
\centering
\includegraphics[width=0.98\columnwidth]{figures/exp/fig_branch_count_ablation.pdf}
\caption{Branch-count ablation on PROTEINS and DD under GCNConv. Accuracy rises
only within an effective operating region before saturating or declining.}
\label{fig:branch_ablation}
\end{figure}

The mechanism metrics explain this split. On PROTEINS, HorizontalRes improves
from \(B=1\) to \(B=4\) while branch diversity rises from 0.0000 to 1.1450 and
branch cosine falls from 1.0000 to 0.8417, but diversity alone is not sufficient:
at \(B=8\) diversity increases further to 1.5441 while accuracy drops to 0.6881.
MatrixRes follows the same pattern at a larger scale, peaking at \(B=6\) with a
diversity of 3.7755. On DD the best settings have nearly saturated branch cosine
(0.9969 for HorizontalRes and 0.9991 for MatrixRes at \(B=2\)) and CKA of
1.0000, so extra residual traffic is unlikely to create new functional
separation; larger branch counts raise diversity without improving accuracy.
MatrixResGated behaves as genuine residual control rather than a dense path: its
active ratio is 0.3121 at the best PROTEINS setting and 0.2352 at the best DD
setting, while the learned gate stays near 0.78--0.84 rather than collapsing.
These are diagnostic correlations rather than causal interventions. The gated
variant does not dominate MatrixRes overall, so sparse filtering is best treated
as a traffic-control mechanism, not a universal improvement.

\subsection{Sensitivity and Five-Fold Validation}
\label{sec:res-sensitivity}

Fold-0 scans only identify candidates. On PROTEINS, MatrixRes does not improve
over its fold-0 baseline in the tested sweep, while MatrixResGated has a
mild-sparsity candidate (\(\lambda=0.02\) raises fold-0 accuracy from 0.6726 to
0.6951; stronger sparsification is harmful). On DD both models prefer lighter
hidden dimensions, reaching 0.7595 and 0.7511 at \(d=32\). Rerunning selected
candidates across all five folds prevents over-claiming: on DD the \(d=32\)
MatrixResGated candidate stays competitive at \(0.7215\pm0.0330\) while cutting
parameters from 42{,}371 to 15{,}043, whereas on PROTEINS the sparsity candidate
does not preserve its fold-0 advantage (\(0.6909\pm0.0447\)) and remains
unstable rather than a confirmed gain.
